% documentation/adaptive_kmesh.tex
%
% Notes on the adaptive k-mesh refinement in LinReTraCe.
% Compile with: pdflatex adaptive_kmesh.tex
%
\documentclass[11pt,a4paper]{article}

\usepackage[margin=2.5cm]{geometry}
\usepackage{amsmath,amssymb}
\usepackage{booktabs}
\usepackage{hyperref}
\usepackage[T1]{fontenc}

\newcommand{\dfda}{\partial f/\partial\varepsilon}
\newcommand{\kk}{\mathbf{k}}
\newcommand{\code}[1]{\texttt{#1}}

\title{Adaptive $k$-mesh refinement in LinReTraCe}
\author{}
\date{}

\begin{document}
\maketitle

\section{Motivation}

Every transport coefficient LinReTraCe evaluates has the structure
\begin{equation}
  L^{ij} \;\propto\; \sum_{\kk n} w_{\kk}\, M^{ij}_{n}(\kk)\,
                     K\!\left(\varepsilon_{n\kk}-\mu\right),
  \label{eq:transport}
\end{equation}
with $M^{ij}$ the optical elements and $K$ a kernel peaked at $\mu$ with
half-width
\begin{equation}
  W \;=\; \max\!\left(k_B T,\ \Gamma\right).
  \label{eq:width}
\end{equation}
For the parameter ranges LinReTraCe is designed for, $W$ can be three or more
orders of magnitude below the bandwidth. A uniform mesh fine enough to place
several distinct band energies inside $W$ is then absurdly dense everywhere
that does not contribute. The adaptive scheme subdivides only where the
integration demonstrably fails.

The entry points are
\begin{verbatim}
  ltb   refine coarse.hdf5 tb_file        <mu> <gamma_min> <T_min> [options]
  lwann refine coarse.hdf5 <w90 folder>   <mu> <gamma_min> <T_min> [options]
\end{verbatim}
where \code{gamma\_min} and \code{T\_min} are the \emph{worst-case corner} of
the sweep to be run afterwards, i.e.\ the values that make $W$ in
Eq.~\eqref{eq:width} smallest.

\section{The error metric}
\label{sec:metric}

\subsection{Sum rule on the energy axis}

Let $B=\{\varepsilon_1<\dots<\varepsilon_N\}$ be the sorted set of distinct
band energies of the current mesh (all bands, all $k$-points). Because
\begin{equation}
  \int_{-\infty}^{\infty} \frac{\partial f}{\partial\varepsilon}\,
  \mathrm{d}\varepsilon \;=\; 1 ,
  \label{eq:sumrule}
\end{equation}
the accuracy with which the trapezoid rule on $B$ reproduces this integral
measures whether the sampled energies resolve the kernel. Writing the
per-panel defect
\begin{equation}
  d_i \;=\; \underbrace{\tfrac{1}{2}\left(\varepsilon_{i+1}-\varepsilon_i\right)
            \left[g(\varepsilon_i)+g(\varepsilon_{i+1})\right]}_{\text{trapezoid}}
       \;-\; \underbrace{\int_{\varepsilon_i}^{\varepsilon_{i+1}}
            g(\varepsilon)\,\mathrm{d}\varepsilon}_{\text{exact}},
  \qquad g \equiv \dfda ,
  \label{eq:defect}
\end{equation}
the metric reported by \code{compute\_error} is
\begin{equation}
  \boxed{\;
  \mathcal{E} \;=\; \sum_i \left|d_i\right|
              \;+\; \left|1-\sum_i \int_{\varepsilon_i}^{\varepsilon_{i+1}}
                    g\,\mathrm{d}\varepsilon\right| \; }
  \label{eq:metric}
\end{equation}
The second term is the truncation of the sampled range $[\varepsilon_1,
\varepsilon_N]$ and is usually negligible ($\sim 2\times10^{-4}$ for
graphene).

The exact panel integrals are evaluated by five-point Gauss--Legendre
quadrature. The analytic kernel is available in closed form, so there is no
reason to estimate $d_i$ by an embedded (bisection) rule; that estimator
collapses precisely where it matters most, on a coarse mesh whose panels are
far wider than the kernel, because neither the trapezoid nor its bisection
resolves the peak. On the graphene $48\times48\times1$ starting mesh it
reported $20.3$ against a true defect of $39.7$. Five nodes is converged: on a
$7.1\times10^{5}$-panel mesh it gives $0.006653$ against $0.006652$ with
twenty nodes, at a fifth of the cost.

\subsection{Why the absolute value is taken per panel}

The historical metric was $\left|\sum_i d_i\right| = \left|1-\int_B
g\right|$, the absolute value of the \emph{signed} sum. That is not a usable
optimisation target. The trapezoid rule underestimates a concave integrand and
overestimates a convex one, and $\dfda$ is a peak: concave over its top,
convex in both tails. The two contributions therefore \emph{always} carry
opposite signs, in every system. Whenever they happen to be of comparable
magnitude, an iterate whose peak defect cancels part of the tail defect scores
better than the strictly finer mesh that follows it. Refinement only ever adds
$k$-points, so a rising signed metric never means a worse mesh --- but the
loop cannot know that.

Graphene, $48\times48\times1$ irreducible start, $\Gamma=1$~meV, $T=1$~K:

\begin{center}
\begin{tabular}{rrr}
\toprule
$N_\varepsilon$ & signed $\left|\sum d_i\right|$ & $\sum\left|d_i\right|$ \\
\midrule
     475 & 0.04480 & 0.04501 \\
     587 & 0.02301 & 0.04242 \\
    1435 & 0.02966 & 0.03135 \\
  714051 & 0.00644 & 0.00665 \\
\bottomrule
\end{tabular}
\end{center}

\noindent
The signed metric dips at $N_\varepsilon=587$ and rebounds at $1435$; the
plateau detector read that as a $-29\%$ regression and stopped. The
per-panel absolute sum falls monotonically throughout.

Since $\left|\sum_i d_i\right|\le\sum_i\left|d_i\right|$, the new metric is
conservative: a mesh that passes it would also have passed the old test, so
nothing previously converged becomes unconverged. The price is that runs which
used to pass through accidental cancellation now need more iterations.

\section{Hotspot selection}
\label{sec:selection}

The loop now optimises the same quantity it reports. Panels are sorted by
$|d_i|$ and marked until their cumulative defect reaches a fraction
\code{--defect\_fraction} (default $0.9$) of the total. A $k$-point is a
hotspot if any of its band energies is an endpoint of a marked panel;
subdividing exactly those $k$-points is what halves the offending panel widths
on the next iteration.

Two earlier design choices were wrong and are worth recording.

\paragraph{Pointwise defects cannot reach the tails.}
The previous criterion flagged energies where the linear interpolant of
$\dfda$ through $B$ deviated pointwise from the analytic kernel, and turned
the spread of those energies into a $\mu$-centred radius. The kernel amplitude
is $1/(\pi\Gamma)$ at $\mu$ against $\Gamma/(\pi\varepsilon^2)$ in the tails, a
ratio of $(\varepsilon/\Gamma)^2$. A pointwise criterion therefore ranks the
peak orders of magnitude above the tails however coarse the tails are. On
graphene the window collapsed onto its floor of $2W$ from the fourth iteration
onward, while the residual sat between $10^{-2}$ and $10^{-1}$~eV; the metric
then fell only as $N^{-0.26}$.

\paragraph{A radius is all-or-nothing.}
Selecting whole panels rather than a radius also avoids dragging in every
$k$-point closer to $\mu$ than the offending one. Forcing the window open to
its $0.1$~eV ceiling did reach the target, but needed $5.2\times10^{5}$ wedge
$k$-points.

\paragraph{Window guard.}
\code{--energy\_window} (default $0.1$~eV) is a \emph{ceiling}, not the
refinement region: it only excludes far-away van Hove defects at the band
edges, which carry real defect but no transport weight. A panel counts as
in-window when it \emph{overlaps} $[\mu-W_{\max},\mu+W_{\max}]$, not when its
centre lies inside --- on a coarse mesh the panel straddling $\mu$ can be far
wider than the window. Graphene $24\times24\times1$ has its first energy above
the Dirac point near $0.25$~eV, so the single panel holding the entire defect
has its centre at $-0.125$~eV; a centre test declared convergence at an error
of $81.9$.

\section{Mesh bookkeeping}

Each hotspot $k$-point is replaced by $n^{d}$ children ($n=$
\code{--refinement\_factor}, odd, default 3; $d$ = number of active axes)
placed at
\begin{equation}
  \kk_\text{child} \;=\; \kk_\text{parent}
    + \left(j-\tfrac{n-1}{2}\right)\frac{\boldsymbol{\delta}_\text{parent}}{n},
  \qquad j=0,\dots,n-1 ,
\end{equation}
where $\boldsymbol{\delta}$ is the per-point cell width stored in
\code{/.kmesh/cell\_deltas}. For odd $n$ the central child coincides with the
parent, so the parent survives every refinement and each iterate is a strict
superset of its predecessor. Weights are split \emph{per parent},
$w_\text{child}=w_\text{parent}/n^{d}$: the parent cell is exactly tiled by
its own children, so this is the only split that conserves weight locally.
Pooling the removed weight and redistributing it uniformly conserves the total
but transfers weight between parents, which is wrong whenever their weights
differ --- the generic case for an irreducible mesh, where the weight encodes
the star multiplicity.

\section{Irreducible meshes}

An irreducible coarse mesh is the cheaper starting point and is fully
supported. Band energies are symmetry invariant, but the optical elements are
not: a raw weighted sum over the wedge breaks the point-group symmetry of the
Onsager tensor. The evaluation therefore replaces the moments at every wedge
point by their group average
\begin{equation}
  M_\text{sym}(\kk) \;=\; \frac{1}{N_\text{sym}}
    \sum_{S} M\!\left(S\kk\right).
\end{equation}
On a regular grid this reproduces the full-zone sum to machine precision for
any lattice.

\paragraph{Two subtleties.}
First, a refined mesh is written with \code{.kmesh/irreducible = False} plus a
\code{.kmesh/symmetrized} marker. The former flag means, to
\code{linretrace}, only that the weights are reconstructible as
$m_{\kk}\,w_\text{sum}/(n_x n_y n_z)$ --- true for a regular Monkhorst--Pack
grid and false for any adaptive mesh, whose weights must be read from
\code{.kmesh/weights}. The marker preserves the symmetrisation for
continuation runs.

Second, on a refined wedge the group average is a valid quadrature but not
identical to refining the reducible mesh, unless the point group acts on the
\emph{fractional} axes by signed permutation. That holds for cubic,
tetragonal, orthorhombic, monoclinic and triclinic lattices; it fails for
hexagonal, trigonal and rhombohedral ones, where a $60^\circ$ rotation carries
the axis-aligned $n\times n$ patch of children onto a sheared one. Both
quadratures carry correct weights and converge to the same integral; they
sample different points. For graphene, $8\times8\times1$ subdivided by 3, the
two differ by $0.11\%$ while each is $0.4\%$ from the converged value.

\section{Stopping}

The loop stops when $\mathcal{E}\le$ \code{--error\_tol} (default
$5\times10^{-3}$), when no panel inside the window still carries a significant
defect, at \code{--max\_iter}, or on a plateau: an iteration improving
$\mathcal{E}$ by less than \code{--plateau\_tol} (default $5\%$), checked from
the fourth refinement step onward. With the metric of
Sec.~\ref{sec:metric} a \emph{rising} $\mathcal{E}$ should no longer occur; if
it does, the loop reports it as a metric failure rather than as a regression,
and keeps the last (finest) mesh.

\section{Performance}

Graphene, $\Gamma=1$~meV, $T=1$~K, default settings, irreducible start:

\begin{center}
\begin{tabular}{lrrr}
\toprule
start & iterations & final $k$-points & $\mathcal{E}$ \\
\midrule
$24\times24\times1$ & 7 &  357 & 0.00112 \\
$48\times48\times1$ & 6 &  425 & 0.00207 \\
\bottomrule
\end{tabular}
\end{center}

\noindent
Under the previous metric and selection the same $48\times48\times1$ run
never converged: it stopped on a false plateau at 1435 $k$-points, and with
the plateau check disabled reached only $\mathcal{E}=0.0064$ at
$3.6\times10^{5}$ $k$-points.

The metric costs one $\dfda$ evaluation per band-axis node plus five per
panel, chunked at $2\times10^{6}$ panels ($\sim80$~MB). Measured at
$7.1\times10^{5}$ panels: $2.6$~s, against minutes for the Hamiltonian
evaluation at that mesh size. On the symmetrising path each iteration
diagonalises $N_k \times N_\text{sym}$ matrices, so an irreducible run trades
a factor $N_\text{sym}$ in work per point against a factor $N_\text{sym}$
fewer points.

\section{Limitation: the metric is a proxy}
\label{sec:proxy}

$\mathcal{E}$ is a pure energy-axis criterion. It never sees the $k$-point
weights, so two meshes with identical energy sets score identically however
differently they sample the Brillouin zone. Resolving the kernel in energy is
necessary for Eq.~\eqref{eq:transport} to be accurate; it is not sufficient.

Concretely: the graphene mesh that reaches $\mathcal{E}=0.0021$ at 425
$k$-points gives a total conductivity about $12\%$ below the converged
dense-grid value, and that deviation is insensitive to
\code{--defect\_fraction} (tested at $0.9$, $0.99$, $0.999$). Treat
\code{--error\_tol} as ``is the kernel resolved in energy'', and confirm
convergence of the observable itself by continuing the refinement or starting
from a finer coarse mesh until it stops moving.

\end{document}
